\chapter{Wykład X.}

\section{Teoria pierścieni.}

Ustalmy zbiór z dwoma działaniami $+$ i $\cdot $: $\left( R, +,  \cdot  \right) $.

Uwaga notacyjna: $x\cdot y + z := \left( x\cdot y \right) + z$ .

Kilka potencjalnych własności tych działań:
\begin{enumerate}[label=A\arabic*.]
	\item $\left( R, + \right) $ jest grupą abelową- element neutralny oznaczamy przez $\mathbf{0}$.
	\item $\left(\forall a,b,c \in R \right)\left( a\cdot \left( b+c \right) = a\cdot b + a\cdot c \land (b+c)\cdot a = b\cdot a + c\cdot a\right) $ -rozdzielność $\cdot $ względem $+$.
	\item $\left(\forall a,b,c\in R \right) \left( a\cdot \left( b\cdot c \right) = \left( a\cdot b \right) \cdot c = a\cdot b\cdot c \right) $- łączność mnożenia.
	\item $\left(\exists \mathbf{1}\in R \right)\left(\forall a\in R \right)\left( a\cdot \mathbf{1} = \mathbf{1}\cdot a = a\right)$- jedynka mnożenia.  ($A4 \text{ i } A3 \implies $ jedyność $\mathbf{1}$).
	\item $\left(\forall a,b\in R \right)\left( a\cdot b = b\cdot a\right)  $-przemienność $\cdot $.
\end{enumerate}
\begin{definition}[Pierścień]
	Dla ustalonego zbioru $R$ oraz działań $+$ i $\cdot $, definiujemy zbiór $\left( R, +, \cdot  \right) $ jako:
	\begin{enumerate}[label=\arabic*.]
		\item \textbf{Pierścień}: jeśli spełnione są warunki $A1, A2$ oraz $A3$.
		\item  \textbf{Pierścień z jednością(z jedynką):} jeśli spełnione są warunki $A1, A2, A3$ oraz $A4$.
		\item \textbf{Pierścień przemienny:} jeśli spełnione są warunki $A1, A2, A3$ oraz $A5$.
		\item \textbf{Pierścień przemienny z jednością(z jedynką):} jeśli spełnione są warunki $A1, A2, A3, A4$ oraz $A5$.
	\end{enumerate}
\end{definition} 

Konwencja (R- pierścień):
\begin{itemize}
	\item $ab:= a\cdot b$,
	\item dla $n\in \N_{>0}$ mamy $a^{n}=\underbrace{a\cdot \ldots\cdot a}_{\text{n razy}}$,
	\item  $a-b := a+ (-b)$
\end{itemize}
\begin{definition}[Podpierścień.]
	Niech $\left( R,+,\cdot  \right) $ będzie pierścieniem. $S\subseteq R$ jest \textbf{podpierścieniem} gdy
	\begin{enumerate}[label=\arabic*.]
		\item $S\le (R,+)$, czyli $S$ podgrupą  $R$,
		\item  $\left(\forall x,y\in S \right)\left( xy \in S \right)  $.
	\end{enumerate}
	\noindent \textbf{Oznaczenie:} $S \underset{\scriptscriptstyle\text{podp}}{\subseteq}  R$ 
\end{definition} 
\begin{remark}
Oczywiście $(S, +\restriction_{S\times S}, \cdot\restriction_{S\times S})$ jest pierścieniem, gdy $S$ jest podpierścieniem $R$.
\end{remark} 

\noindent \textbf{Przykłady:} 
\begin{enumerate}[label=\arabic*)]
	\item $\left( \mathbb{Z} , +, \cdot  \right) $ : pierścień przemienny z $\mathbf{1}$.
	\item $\left( \N, +, \cdot  \right) $ : \textbf{nie} jest pierścieniem. 
	\item $\left( 2\mathbb{Z}, +, \cdot  \right) $: pierścień przemienny (bez $\mathbf{1}$ ).
	\item $\left( \mathbb{Z} _{n}, +_{\scriptscriptstyle n}, \cdot _{\scriptscriptstyle n} \right) $.
	\item $\left( \mathbb{Q} , +,\cdot  \right) $, $\left( \mathbb{R} , +, \cdot  \right) $, $\left( \mathbb{C}, +, \cdot  \right) $ : pierścienie przemienne z $\mathbf{1}$.

	Ponadto \[
	2\mathbb{Z} \underset{\scriptscriptstyle\text{podp}}{\subseteq} \mathbb{Z} \underset{\scriptscriptstyle\text{podp}}{\subseteq} \mathbb{Q} \underset{\scriptscriptstyle\text{podp}}{\subseteq} \mathbb{R} \underset{\scriptscriptstyle\text{podp}}{\subseteq} \mathbb{C} 
	.\] 
	\item $\left( M_{n}\left( \mathbb{R}  \right), +, \cdot  \right) $ : pierścień z $\mathbf{1}$.

	Jest on przemienny $\iff n=1$.
    	\item Niech R będzie pierścieniem przemiennym z $\mathbf{1}$. Tak samo jak w $M_{n}\left( \mathbb{R}  \right) $ definiujemy dodawanie i mnożenie macierzy (używamy $+$ i $\cdot $ z $R$ ).
	
	Wtedy $\left( M_{n}\left( \mathbb{R}  \right)  \right) $ to pierścień z $\mathbf{1}$, tzn.
	\[
		\mathbf{1} = \mathcal{I} = \begin{pmatrix} \mathbf{1}_{R} & \ldots & 0 \\ \vdots & \ddots & \vdots \\ 0 & \ldots & \mathbf{1}_{R} \end{pmatrix} 
	.\] 
	Czyli przykładowo mamy pierścień $M_3\left( \mathbb{Z}_{3} \right) $.
	\item Niech $X$- zbiór oraz $R$- pierścień, wtedy $\underbrace{\text{Fun}\left( X,R \right)}_{R^{X}}$ -pierścień z działaniem 
		\begin{align*}
			\left( f + g \right) \left( x \right) &:= f\left( x \right) +_{R} g\left( x \right), \\
			\left( f \cdot g \right) &:= f\left( x \right) \cdot_{R} g\left( x \right) 
		.\end{align*}
		oraz
		\begin{align*}
			R: \text{ pierścień przemienny } &\implies \text{Fun}\left( X,R \right) \text{ przemienny},\\
			R: \text{ pierścień przemienny z} \mathbf{1} &\implies \text{Fun}\left( X,R \right) \text{ przemienny z } \mathbf{1}.
		.\end{align*}
	\item Definiujemy podpierścienie $\text{Fun}\left( R, R \right)$:
		\begin{itemize}
			\item $C\left( R \right) $ : pierścień funkcji ciągłych.
			\item $C^{n}\left( R \right)$: pierścień funkcji $n$-krotnie różniczkowalnych,
			\item $C^{\infty}$ : pierścień funkcji nieskończenie wiele różniczkowalnych (gładkich).
		\end{itemize}
	\item Niech $A$ będzie grupą abelową. Wtedy
		\[
		\text{End}\left( A \right) := \{f: A\ to A: f \text{ jest endomorfizmem grupy }\left( A,+ \right) \}
		.\] 
		Dla $f,g \in \text{End}\left( A \right) $ definiujemy:
		\begin{enumerate}[label=\roman*)]
			\item $f+g: A \to A$ wzorem
				\[
					\left( f+g \right) \left( a \right) := f\left( a \right) +_{A} g\left( a \right) 
				.\]
			(WYMAGA KOMENTARZA czy to nie wynika samo przez się?)
		        \item $f\cdot g: A \to A$ wzorem
				\[
					\left( f\cdot g \right)\left( a \right) := f\left( a \right) \cdot_{A} g\left( a \right) 
				.\] 
		\end{enumerate}
		Wtedy $\left( \text{End}\left( A \right), +, \cdot  \right) $ jest pierścieniem z $\mathbf{1}=\id_{A}$, zwykle nieprzemienny.
\end{enumerate}
\begin{fact}
	Niech 
	\begin{itemize}
		\item $R$ będzie pierścieniem,
		\item $a,b,c \in R$,
		\item $n,m \in \N_{>0}$,
		\item $\alpha , \beta \in \mathbb{Z} $.
	\end{itemize}
	Wtedy
	\begin{enumerate}[label=\arabic*)]
		\item $a^{n}b^{m} = a^{n+m}$,
		\item $\left( a^{n} \right)^{m} = a^{nm}$,
		\item $\alpha a + \beta a = \left( \alpha +\beta  \right) a$,
		\item $\alpha \left( \beta a \right) = \left( \alpha \beta  \right) a$,
		\item $a\mathbf{0} = \mathbf{0} = \mathbf{0}a$,
		\item $\left( -a \right) b = -\left( ab \right) = a\left( -b \right) $,
		\item $a\cdot \left(b-c  \right) = ab-ac$,
		\item $\left( b-c \right) a = ba-ca$.
		\item $R$-przemienny $\implies \left( a+b \right)^{n} = \sum_{n = 1}^{\infty} \begin{pmatrix} n \\ i \end{pmatrix} a^{i}b^{n-1}$,
		\item $R$ z $\mathbf{1} \implies \left( -\mathbf{1} \right) a = -a = a\left( -\mathbf{1} \right) $.
	\end{enumerate}
\end{fact}
\begin{proof}
	TODO
\end{proof} 
\begin{remark}
	Z faktu $5)$ mamy: 
	\[
		\left( R, \cdot  \right) \text{ jest grupą }\iff R = \underbrace{ \{\mathbf{0}\} }_{\text{pierścień zerowy}}\iff \mathbf{0} = \mathbf{1}
	.\] 
\end{remark} 
\begin{remark}
	$A2$ mówi dokładnie, że dla  $\left(\forall r\in R \right) $ funkcja
	\[
	\lambda_{r}: R \to R\text{, dana wzorem } \lambda_{r}\left( x \right) := r\cdot x
	.\] 
	jest \textbf{endomorfizmem} grupy $\left( R,+ \right) $.
\end{remark} 
\begin{remark}
	Od teraz  $\mathbb{Z} _{n}, \mathbb{Z} ,\mathbb{Q} ,\mathbb{R} ,\mathbb{C} $ będą oznaczać pierścienie.
\end{remark} 

\begin{definition}[Element odwracalny pierścienia]
	Niech $R$ będzie pierścieniem z $\mathbf{1}$ oraz $a\in R$.
	\begin{enumerate}[label=\arabic*)]
		\item $a$ jest \textbf{odwracalny} gdy
			\[
			\left(\exists b\in R \right)\left( ab = ba = \mathbf{1} \right)  
			.\] 
		\item $R^{\ast }:= \{a\in R: a \text{ jest odwracalny}\} $
	\end{enumerate}
\end{definition} 
\begin{fact}
      $\left( R^{\ast}, \cdot \right) $ jest grupą.
\end{fact}
\begin{proof}
	TODO
\end{proof} 
\begin{remark}
	$R \neq \{\mathbf{0}\} \implies \mathbf{0}\not\in R^{\ast}$
\end{remark} 
\begin{definition}
	Niech $R$ będzie pierścieniem z $\mathbf{1}$. Wtedy
	\begin{enumerate}[label=\arabic*)]
		\item $\left( R,+ \right) $ jest \textbf{grupą addytywną} $R$,
		\item $\left( R^{\ast}, \cdot  \right) $ jest \textbf{grupą multiplikatywną} $R$.
	\end{enumerate}
\end{definition} 
\begin{remark}
	Niech $R$ będzie pierścieniem z $\mathbf{1}$.
	Wtedy
	\[
	\left(\forall a \in R^{\ast} \right)\left(\exists !b\in R^{\ast} \right)\left( ab = ba = \mathbf{1} \right)   
	.\] 
	Oznaczamy wtedy $b = a^{-1}$.

	\noindent \textbf{Ogólniej:}

	$\left(\forall n\in \mathbb{Z}  \right)\left(\forall a \in R^{\ast} \right)  $ mamy zdefiniowane $a^{n}$. Wtedy:
	\begin{itemize}
		\item $a^{-1}$ jest elementem odwrotnym do $a$ w sensie mnożenia,
		\item $-a$ jest elementem odwrotnym do $a$ w sensie dodawania (czyli elementem \textbf{przeciwnym} do $a$ ).
	\end{itemize}
\end{remark}
\begin{remark}
	Dodawanie ($+$ ) nie jest działaniem na $R^{\ast}$ (chyba że $R=\{\mathbf{0}\} $ ).
\end{remark} 

\noindent \textbf{Przykłady:} 
\begin{enumerate}[label=\arabic*)]
	\item $\mathbb{Z} ^{\ast}=\{-1, 1\} $,
	\item $\left( \mathbb{Z} _{n} \right)^{\ast} = \mathbb{Z}_{n}^{\ast} = \{a \in \mathbb{Z} _{n}: \text{NWD}\left( a,n \right) = 1\} $ ($\cong \aut\left(\mathbb{Z}_{n}\right) $),
	\item $\mathbb{Q} ^{\ast} = Q\setminus \{\mathbf{0}\} $, $\mathbb{R} ^{\ast} = \mathbb{R} \setminus \{\mathbf{0}\}$, $\mathbb{C} ^{\ast} = \mathbb{C} \setminus \{\mathbf{0}\} $.
	\item $M_{n}\left( \mathbb{R}  \right) ^{\ast} = GL_{n}\left( \mathbb{R}  \right) $. Podobnie dla $\mathbb{Q} \text{ oraz }\mathbb{C} $.
	\item Niech $R$ będzie pierścieniem przemiennym z $\mathbf{1}$.
		Mamy 
		\[
		GL_{n} (R) := \{A\in M_{n}\left( R \right) : \det\left( A \right) \in R^{\ast} \} 
		.\] 
		Wtedy
		\[
		GL_{n} \left( R \right) = M_{n}\left( R \right) ^{\ast} 
		.\] 
	\noindent Przykładowo
	\begin{align*}
		GL_{n} \left( \mathbb{Z}  \right) &= \{A\in M_{n} \left( \mathbb{Z}  \right): \det\left( A \right) \in \{-1,1\}  \}, \\
		GL_{n} \left( \mathbb{Z} _{n}  \right) &= \{A\in M_{n} \left( \mathbb{Z}  \right): \det\left( A \right) \in \{1,3\}  \}, \\
		SL_{n} \left( R \right) &= \{A \in GL_{n} \left( R \right) : \det\left( A \right)  = 1\} 
	.\end{align*}
	W trzeciej równości i.e. dotyczącej $SL_{n} \left( R \right) $ mamy $SL_{n} \left( R \right) \unlhd GL_{n} \left( R \right) $ i dostajemy 
	\[
	DSL_{n}\left( R \right) := \mfaktor{SL_{n} \left(R\right)}{Z\left( SL_{n}\left( R \right)   \right) } 
	.\] 
	\item Niech $X$ będzie zbiorem oraz $R$ pierścieniem z $\mathbf{1}$. Wtedy
		\[
		\text{Fun}\left( X,R \right) ^{\ast} = \text{Fun}\left( X, R^{\ast}  \right) 
		.\] 
	\item Niech $A$ będzie grupą abelową, wtedy
		\[
		\text{End}\left( A \right) ^{\ast} = \aut\left(A\right) 
		.\] 
\end{enumerate}
\section{Teoria podzielności.}
\begin{context}

	Zakładamy teraz, że $R$ jest \textbf{pierścieniem przemiennym}.
\end{context}
\begin{definition}
	Niech $x,y \in R$. 
	\begin{enumerate}[label=\arabic*)]
		\item $x \mid y$, gdy $\left(\exists r\in R \right)\left( y = rx \right) $ 
		\item $x\sim y$ ($x$ jest \textbf{stowarzyszony} z $y$, gdy $x \mid y$ oraz $y \mid x$
	\end{enumerate}
\end{definition} 
\begin{eg}
	Dla $x,y \in \mathbb{Z} $ 
	\[
	x\sim y \iff x = \pm y 
	.\] 
\end{eg} 
\begin{fact}

	Niech $x,y,z \in R$.
	\begin{enumerate}[label=\arabic*)]
		\item $\left(\exists n \in R^{\ast}\right) \left( y = nx \right)\implies y\sim x  $,
		\item $x \mid y$ oraz $y \mid z$ $\implies x \mid z $,
		\item $R$ jest z $\mathbf{1}$ $\implies \sim  $ jest relacją równoważności.
	\end{enumerate}
\end{fact}
\begin{proof}
	TODO
\end{proof} 
\begin{definition}[Dzielnik zera.]
	$x$ jest dzielnikiem zera, gdy:
	\begin{enumerate}[label=\arabic*)]
		\item $x\neq \mathbf{0}$,
		\item $\left(\exists y\neq \mathbf{0} \right)\left( xy =0 \right)  $.
	\end{enumerate}
\end{definition} 
\begin{definition}[Dziedzina]
	$R$ jest dziedziną (całkowitości), gdy:
	\begin{enumerate}[label=\arabic*)]
		\item $R$ jest pierścieniem przemiennym z $\mathbf{1}$.
		\item $\mathbf{0}\neq \mathbf{1}$ (tzn. $R \neq \{\mathbf{0}\}) $.
		\item $R$ nie ma dzielnika $\mathbf{0}$.
	\end{enumerate}
\end{definition} 
\begin{fact}
	Niech $R$ będzie pierścieniem przemiennym z $\mathbf{1}$ oraz $x,y \in R$. Wtedy:
	\begin{enumerate}[label=\arabic*)]
		\item $R$ jest dziedziną $\implies$ 
	\end{enumerate}
\end{fact}

\noindent \textbf{Przykłady:}
\begin{enumerate}[label=\arabic*)]
	\item $\mathbb{Z}, \mathbb{Q} , \mathbb{R} , \mathbb{C} $ są dziedzinami.
	\item $2\in \mathbb{Z} _{4}$ jest dzielnikiem $\mathbf{0}$ (bo $2\neq 0$ i $2-2=0$).

		Czyli $\mathbb{Z}_{4}$ nie jest dziedziną. 
	\item Podpierścień z $\mathbf{1}$ dziedziny jest dziedziną.
\end{enumerate}
\begin{fact}
	Niech $R$ będzie pierścieniem przemiennym z $\mathbf{1}$ oraz $x,y \in R$. Wtedy
	\begin{enumerate}[label=\arabic*)]
		\item $R$ jest dziedziną $\implies \left( x\sim y \iff \left(\exists n \in R^{\ast}  \right)\left( y = nx \right)   \right) $.
		\item $x\in R^{\ast} \implies x$ nie jest dzielnikiem $\mathbf{0}$.
	\end{enumerate}
\end{fact}
\begin{proof}
	TODO
\end{proof} 
\begin{remark}
	Zauważmy, że
	\begin{enumerate}[label=\arabic*)]
		\item Istnieją pierścienie z $\mathbf{1}$, gdzie $a\sim b$ ale nie ma $a\in R^{\ast} $, takiego że $a=ab$ (być może  $C\left[ 0,1 \right] $?)
		\item W dowodzie pojawiło się \textbf{prawo skracania dla dziedzin}:

			Dla $R$ będącym dziedziną, oraz $a,b \in R$ 
			\[
			ab = ac\text{ oraz }a\neq \mathbf{0}\implies b = c
			.\] 
	\end{enumerate}
\end{remark} 

\begin{definition}[Pierścień z dzieleniem i Ciało]
    Niech $R$ będzie pierścieniem z jednością $\mathbf{1}$.
    \begin{enumerate}[label=\arabic*)]
        \item $R$ nazywamy \textbf{pierścieniem z dzieleniem} (lub pierścieniem z dzielnikiem), gdy każdy jego niezerowy element jest odwracalny, czyli:
        \[ R^{\ast} = R \setminus \{\mathbf{0}\} \]
        Warunek ten implikuje, że pierścień jest nietrywialny ($\mathbf{0} \neq \mathbf{1}$).
        \item $R$ nazywamy \textbf{ciałem}, gdy jest przemiennym pierścieniem z dzieleniem.
    \end{enumerate}
\end{definition}


\textbf{Przykłady:}
\begin{enumerate}[label=\arabic*)]
	\item $\mathbb{Q}, \mathbb{R} , \mathbb{C}  $ są ciałami.
	\item $\mathbb{Z} $ nie jest ciałem.
	\item Niech $n\in \N_{>0}$. Wtedy następujące warunki są równoważne:
		\begin{enumerate}[label=\roman*)]
			\item $\mathbb{Z}_{n}$.
			\item $\mathbb{Z} _{n} $ jest dziedziną.
			\item $n$ jest liczbą pierwszą.
		\end{enumerate}
		\noindent Ogólniej: $i)\iff ii)$ dla pierścieni przemiennych skończonych (dowód na liście zadań).
\end{enumerate}
\noindent \textbf{Oznaczenie: } $\mathbb{F}_{p}: \mathbb{Z}_{p}$ ciało p-elementowe. 

\begin{remark}
	Mamy:
	ciało $\implies$ dziedzina $\implies$ pierścień przemienny z $\mathbf{1}$ $\implies$ pierścień z $\mathbf{1}$.

	Ale
	\[
	\mathbb{R} \nLeftarrow \mathbb{Z} \nLeftarrow \mathbb{Z} _{n} \nLeftarrow M_{2}\left( \mathbb{R}  \right) 
	.\] 
\end{remark} 
% START

\noindent \textbf{Przykłady i klasyfikacja struktur:}
\begin{enumerate}[label=\arabic*)]
    \item $\mathbb{Q}, \mathbb{R}, \mathbb{C}$ są ciałami[cite: 58, 85].
    \item $\mathbb{Z}$ jest dziedziną, ale \textbf{nie} jest ciałem, ponieważ dla $a \in \mathbb{Z} \setminus \{-1, 1\}$ nie istnieje element odwrotny wewnątrz $\mathbb{Z}$[cite: 58, 123].
    
    \item \textbf{Kwaterniony ($\mathbb{H}$):} Jest to pierścień z dzieleniem (skew field), który \textbf{nie} jest ciałem. 
    \begin{itemize}
        \item Każdy niezerowy kwaternion $q = a + bi + cj + dk$ posiada element odwrotny $q^{-1} = \frac{\bar{q}}{|q|^2}$, gdzie $\bar{q} = a - bi - cj - dk$.
        \item Mnożenie nie jest przemienne: zachodzą relacje $ij = k$ oraz $ji = -k$, co wyklucza $\mathbb{H}$ z definicji ciała.
    \end{itemize}

    \item Niech $n \in \mathbb{N}_{>0}$. Następujące warunki są równoważne[cite: 59]:
    \begin{enumerate}[label=\roman*)]
        \item $\mathbb{Z}_n$ jest ciałem[cite: 59].
        \item $\mathbb{Z}_n$ jest dziedziną[cite: 59].
        \item $n$ jest liczbą pierwszą[cite: 59].
    \end{enumerate}
    \textit{Ważne:} Równoważność $i) \iff ii)$ zachodzi dla wszystkich pierścieni przemiennych \textbf{skończonych}:
    \begin{theorem}
Każda skończona dziedzina całkowitości jest ciałem.
\end{theorem}

\begin{proof}[Szkic dowodu]
Niech $R = \{x_1, x_2, \dots, x_n\}$ będzie skończoną dziedziną całkowitości. 
Aby wykazać, że $R$ jest ciałem, musimy udowodnić, że każdy element $a \in R \setminus \{0\}$ posiada odwrotność multiplikatywną.

\begin{enumerate}
    \item \textbf{Konstrukcja odwzorowania:} Dla ustalonego $a \neq 0$ rozważmy odwzorowanie mnożenia lewostronnego $\lambda_a: R \to R$ dane wzorem:
    \[ \lambda_a(x) = ax \]
    
    \item \textbf{Iniekcja (Różnowartościowość):} Załóżmy, że $\lambda_a(x) = \lambda_a(y)$. Wtedy $ax = ay$, co implikuje $a(x-y) = 0$. 
    Ponieważ $R$ jest dziedziną i $a \neq 0$, nie istnieją dzielniki zera, więc musi zachodzić $x-y = 0$, czyli $x=y$. 
    Odwzorowanie $\lambda_a$ jest zatem iniekcją.
    
    \item \textbf{Surjekcja i Jedynka:} Ponieważ $R$ jest zbiorem skończonym, każda iniekcja ze zbioru w ten sam zbiór jest automatycznie surjekcją (zasada szufladkowa Dirichleta). 
    Skoro $\lambda_a$ jest surjekcją, to w obrazie tego odwzorowania musi znajdować się element neutralny mnożenia $\mathbf{1} \in R$.
    
    \item \textbf{Konkluzja:} Istnieje zatem taki element $b \in R$, że $\lambda_a(b) = \mathbf{1}$, co oznacza $ab = \mathbf{1}$. 
    Z przemienności pierścienia wynika, że $b$ jest szukaną odwrotnością elementu $a$.
\end{enumerate}
Zatem każdy niezerowy element $a$ jest odwracalny, co czyni $R$ ciałem.
\end{proof}
	
\end{enumerate}

\noindent \textbf{Wyjaśnienie oznaczenia $\mathbb{F}_p$:} 
W literaturze symbol $\mathbb{F}_p$ (od ang. \textit{field}) oznacza abstrakcyjne ciało o charakterystyce $p$ mające $p$ elementów[cite: 60]. W Twoich notatkach jest ono utożsamiane z $\mathbb{Z}_p$, ponieważ zbiór reszt modulo $p$ (gdzie $p$ jest liczbą pierwszą) dostarcza nam konkretnej, konstruktywnej realizacji takiego ciała.



\begin{remark}[Hierarchia struktur]
Mamy następujący ciąg inkluzji (każda struktura po prawej jest ogólniejsza)[cite: 61, 122]:
\[ \text{Ciało} \implies \text{Dziedzina} \implies \text{Pierścień przemienny z } \mathbf{1} \implies \text{Pierścień z } \mathbf{1} \]

Implikacje w drugą stronę (z prawa na lewo) zazwyczaj nie zachodzą, co pokazują poniższe kontrprzykłady:
\begin{itemize}
    \item \textbf{Pierścień z $\mathbf{1} \nLeftarrow \mathbb{R}$}: Pierścień macierzy $M_n(\mathbb{R})$ dla $n > 1$ posiada $\mathbf{1}$, ale nie jest przemienny[cite: 21, 89].
    \item \textbf{Pierścień przemienny z $\mathbf{1} \nLeftarrow \mathbb{Z}_n$}: Pierścień $\mathbb{Z}_n$ dla $n$ złożonego (np. $\mathbb{Z}_4$) jest przemienny, ale nie jest dziedziną, bo posiada dzielniki zera ($2 \cdot 2 \equiv 0 \pmod 4$)[cite: 53, 115].
    \item \textbf{Dziedzina $\nLeftarrow \mathbb{Z}$}: Liczby całkowite $\mathbb{Z}$ są dziedziną, ale nie są ciałem, bo tylko $1$ i $-1$ są odwracalne[cite: 109, 123].
\end{itemize}
\end{remark}

% FINAŁ

\begin{definition}[Podciało.]
	Niech $L$ będzie ciałem oraz $K\subseteq L$ podzbiorem $L$. Wtedy $K$ jest \textbf{podciałem} $L$, gdy:
	\begin{enumerate}[label=\arabic*)]
		\item $K\underset{\scriptscriptstyle\text{podp}}{\subseteq} L$,
		\item $\mathbf{1} \in K$,
		\item $\left(\forall a \in K\setminus \{\mathbf{0}\}  \right) \left( a^{-1}\in K \right) $
	\end{enumerate}
\noindent Oznaczamy $K\underset{\scriptscriptstyle\text{podc}}{\subseteq} L$ .
\end{definition} 
\begin{remark}
	Wtedy $\left( K, +\restriction_{\scriptscriptstyle K\times K}, \cdot\restriction_{\scriptscriptstyle K\times K} \right) $ jest ciałem.
\end{remark} 
\noindent \textbf{Przykłady:}
\begin{enumerate}[label=\arabic*)]
	\item $\mathbb{Q} \underset{\scriptscriptstyle\text{podc}}{\subseteq} \mathbb{R} \underset{\scriptscriptstyle\text{podc}}{\subseteq} \mathbb{C} $.
	\item $\mathbb{Z} \underset{\scriptscriptstyle\text{podp}}{\subseteq} \mathbb{Q} $, ale  $\mathbb{Z} $ nie jest podciałem $\mathbb{Q} $.
\end{enumerate}
